\documentclass[aspectratio=169]{beamer}
\usetheme{Madrid}
\usecolortheme{default}
\usepackage{amsmath}
\usepackage{graphicx}
\usepackage{tikz-cd}
\usepackage{multicol}

\setlength{\parindent}{0pt}
\setlength{\parskip}{1\baselineskip}

\title[Lecture 02]{An Introduction to E-Commerce\\\small with a Focus on Machine Learning, Deep Learning, and Artificial Intelligence}
\subtitle{Lecture 02: Platforms, Marketplaces, and Business Models}
\author{Dr. Haitham A. El-Ghareeb}
\institute{Faculty of Computers and Information Sciences \\ Mansoura University \\ Egypt}
\date{\today}

\begin{document}

\begin{frame}
  \titlepage
\end{frame}

\begin{frame}{Session plan (90 minutes)}
  \begin{enumerate}
    \item From storefront to platform
    \item Value creation, value capture, and growth
    \item Governance and trust as system requirements
    \item Business model patterns (and architectural implications)
    \item Where AI fits (system-level view)
    \item Hands-on (placeholder)
  \end{enumerate}
\end{frame}

\begin{frame}{Learning objectives}
  By the end of this lecture, you should be able to:
  \begin{itemize}
    \item Distinguish e-commerce \emph{channels} (storefronts) from \emph{platforms} (ecosystems) and \emph{marketplaces} (multi-seller).
    \item Analyze platform economics: network effects, multi-homing, and pricing of participation.
    \item Translate business model choices into architectural requirements (identity, trust, data, integration, ERP).
  \end{itemize}
\end{frame}

\section{From storefront to platform}

\begin{frame}{From storefront to platform}
  \begin{itemize}
    \item \textbf{Storefront:} a single merchant selling to buyers (B2C/B2B) via a web/app channel.
    \item \textbf{Platform:} a system that enables interactions among multiple participant groups (e.g., buyers, sellers, advertisers, logistics providers).
    \item \textbf{Marketplace:} a platform where multiple sellers list items and transactions are mediated by platform policies and infrastructure.
  \end{itemize}
\end{frame}

\section{Value creation, value capture, and growth}

\begin{frame}{Value creation, value capture, and growth}
  \begin{itemize}
    \item \textbf{Value creation and matching:} Platforms create value by reducing search and transaction costs, increasing variety, and improving trust; AI frequently strengthens these effects by improving matching and decision support.
    \item \textbf{Value capture models:} Take rate (commission on completed transactions); Subscription (seller or buyer membership tiers)
    \item \textbf{Network effects:} \textbf{Direct network effects:} more users attract more users (e.g., social commerce communities).; \textbf{Indirect network effects:} more buyers attract more sellers and vice versa (classic marketplace dynamic).
  \end{itemize}
\end{frame}

\section{Governance and trust as system requirements}

\begin{frame}{Governance and trust as system requirements}
  \begin{itemize}
    \item Seller onboarding, verification, and compliance workflows (KYC/KYB where relevant).
    \item Catalog integrity: listing policies, moderation, and counterfeit detection.
    \item Reviews/ratings, dispute resolution, refunds/returns policy enforcement.
    \item Trust \& safety operations: rule+ML detection, case management, and human-in-the-loop escalation.
  \end{itemize}
\end{frame}

\section{Business model patterns (and architectural implications)}

\begin{frame}{Business model patterns (and architectural implications)}
  \begin{itemize}
    \item \textbf{B2C retail:} Key systems: PIM/catalog, promotions, payments, fulfillment, customer support.; ERP/CRM integration: order-to-cash, inventory, accounting, returns.
    \item \textbf{B2B commerce:} Quotes, negotiated pricing, approval workflows, invoicing, credit limits.; Integration-heavy: procurement, supplier management, and EDI/API flows.
    \item \textbf{Marketplace:} Split catalog ownership (platform vs sellers), seller tools, payout/settlement.; Additional "control plane": seller policies, listing moderation, abuse prevention.
    \item \textbf{Digital services (subscriptions):} Recurring billing, entitlements, and churn/retention analytics.; Usage-based pricing requires event design and metering.
  \end{itemize}
\end{frame}

\section{Where AI fits (system-level view)}

\begin{frame}{Where AI fits (system-level view)}
  \begin{itemize}
    \item \textbf{Matching layer:} search, recommendations, ranking.
    \item \textbf{Trust layer:} fraud/abuse detection, content moderation.
    \item \textbf{Operations layer:} forecasting, replenishment, routing.
    \item \textbf{Experience layer:} personalization and conversational commerce.
  \end{itemize}
\end{frame}

\section{Hands-on (placeholder)}

\begin{frame}{Hands-on (placeholder)}
  \begin{itemize}
    \item Create a one-page \emph{platform blueprint}: participant groups, value exchange, pricing model, and required capabilities.
    \item Map capabilities to an EA view and identify which capabilities are ERP-owned vs commerce-owned.
  \end{itemize}
\end{frame}

\end{document}
